\fontfam[LM]
\load[emoji]
\typosize[10/12]
\hyperlinks \Blue \Green

\centerline{\typosize[20/24]\bf State-space exploration}
\smallskip
\centerline{\typosize[14/18]\bf Search algorithms for Clojure}
\medskip
\centerline{A~lazy Clojure library for implicit state spaces}
\smallskip
\centerline{\it Matteo Caoduro}
\bigskip
\nonum
\sec Contents
\maketoc

\sec User documentation

State-space is a~Clojure library for lazily exploring graphs and generated
state spaces. Graph searches yield visits lazily: maps whose \code{:state}
field contains a~value from the problem domain, with optional information
about how that state was reached. Examples use \code{vec} to show all visits
for small searches.

To try the examples from the project directory, start a~REPL with
\code{clojure -M}; the source directory is on the classpath. Load the library
with:

\begtt \hisyntax{clojure}
(require '[state-space.core :as search])
\endtt

The algorithm sections introduce the core searches in order. The later
sections on visit identity, additional visit data, and termination are
reference material; the worked problems show where those options are useful.
For a~first pass, focus on the algorithm descriptions and then follow the
worked problems in section~\ref[examples]{Worked problems}.

\secc Search spaces and depth-first search

A~state is a~value from the problem domain, such as a~location, puzzle board,
or configuration. A~{\it space} says which states can follow a~state. A~space
may be an adjacency map or a~function that returns successors for a~state.
DFS explores a~visit's successors before returning to pending siblings.
Starts and successors are considered in the order supplied.

The simplest space is an adjacency map. Maps are callable in Clojure, so
they can be passed directly:

\begtt \hisyntax{clojure}
(def graph
  {:start [:left :right]
   :left  [:goal]
   :right [:goal]
   :goal  []})

(vec (search/dfs graph [:start]))
;; => [{:state :start} {:state :left}
;;     {:state :goal} {:state :right}]
\endtt

The first route reaches \code{:goal} through \code{:left}, so the default
behavior reports that state once and does not report its later arrival
through \code{:right}. A~missing map entry is also fine: it has no
successors and is treated as a~leaf.

A~space can be a~function instead of a~stored graph. By default, it receives
the current \code{:state} and returns a~collection of successor states:

\begtt \hisyntax{clojure}
(defn successors [n]
  (when (< n 3)
    [(inc n)]))

(mapv :state (search/dfs successors [0]))
;; => [0 1 2 3]
\endtt

DFS is lazy: it does not construct the whole state space before returning
visits. The consumer can take a~finite prefix even when the space is infinite:

\begtt \hisyntax{clojure}
(mapv :state
      (take 4 (search/dfs (fn [n] [(inc n)]) [0])))
;; => [0 1 2 3]
\endtt

Successors of a~yielded visit are requested only when the consumer asks for
more results.

\secc[context] Visits and route context

The simplest visit contains just \code{:state}. When a~search needs to keep
track of how it reached a~state, extensions add context fields such as
\code{:depth} or \code{:path}. An extension is a~map of instructions for the
search, not data copied into each visit. The search calls its
\code{:start} and \code{:next} hooks, then merges the map each returns into
the visit. For example, the depth extension behaves like this:

\begtt \hisyntax{clojure}
{:start (fn [_visit] {:depth 0})
 :next  (fn [parent _transition _visit]
          {:depth (inc (:depth parent))})}
\endtt

\code{:start} initializes each starting visit, and \code{:next} computes
context for a~newly reached visit from its parent. The hook names say when
the callbacks run; keys such as \code{:depth} are the context they return.
Post-order traversal's \code{:exit} hook also returns context, while its
\code{:more?} hook returns a~boolean to control whether another successor
is explored. The built-in depth extension packages the common
\code{:start}/\code{:next} pattern, so callers can simply pass it as an
option:

\begtt \hisyntax{clojure}
(mapv (juxt :state :depth)
      (search/dfs graph [:start]
                  {:extensions [search/depth-extension]}))
;; => [[:start 0] [:left 1] [:goal 2] [:right 1]]
\endtt

The path extension records the states on the route so far:

\begtt \hisyntax{clojure}
(mapv :path
      (search/dfs graph [:start]
                  {:extensions [search/path-extension]}))
;; => [[:start] [:start :left] [:start :left :goal]
;;     [:start :right]]
\endtt

The path extension uses the same mechanism. Extensions run in order; maps
returned by later callbacks override same-named values from earlier ones.
Post-order traversal also supports \code{:exit} and \code{:more?} hooks,
described with its other callbacks in the additional-data reference section.

\secc Breadth-first search

BFS visits all states at one edge-count depth before visiting states at the
next depth. It is useful when the goal is a~state nearest to any of the
starting states in an unweighted space. For example, DFS follows the detour
first in this graph, while BFS finds the nearer route:

\begtt \hisyntax{clojure}
(def routes {:start [:detour :near]
             :detour [:middle]
             :middle [:goal]
             :near [:goal]
             :goal []})

(vec (search/bfs routes [:start]))
;; => [{:state :start} {:state :detour} {:state :near}
;;     {:state :middle} {:state :goal}]
\endtt

By default, BFS reports each state only on its first arrival. In a~static
space, this means it reports a~state at its minimum edge-count depth from any
start; here the first route to \code{:goal} is the two-edge route through
\code{:near}. See section~\ref[identity]{Visit identity} for how to change
which arrivals count as distinct.
If multiple routes reach a~state at that depth, start order and successor
order determine which route is retained. BFS does not account for edge
weights; use it when every transition has equal cost.

The basic search needs no extensions; each visit contains its state. Add the
depth extension when the edge-count depth is useful:

\begtt \hisyntax{clojure}
(mapv (juxt :state :depth)
      (search/bfs routes [:start]
                  {:extensions [search/depth-extension]}))
;; => [[:start 0] [:detour 1] [:near 1]
;;     [:middle 2] [:goal 2]]
\endtt

To retrieve the shortest route, add the path extension, then use the standard
Clojure combination \code{first} and \code{filter} to take the first matching
visit. Since the search is lazy, it stops as soon as the goal is found:

\begtt \hisyntax{clojure}
(->> (search/bfs routes [:start]
                 {:extensions [search/path-extension]})
     (filter #(= :goal (:state %)))
     first
     :path)
;; => [:start :near :goal]
\endtt

All starting visits are at depth zero and are visited before any successors.
BFS is lazy too. A~bounded prefix can be consumed even from an infinite
space:

\begtt \hisyntax{clojure}
(mapv :state
      (take 4 (search/bfs (fn [n] [(inc n)]) [0])))
;; => [0 1 2 3]
\endtt

\secc Simple paths

\code{simple-paths} enumerates routes depth-first without revisiting a~state
within a~route (by default). Unlike DFS, it does not discard a~state merely
because a~different route reached it first. This is useful when all
cycle-free ways through a~small graph matter. Like the other searches, it
first returns visits with just their \code{:state}. Here the repeated
\code{:c} is reached by two distinct routes:

\begtt \hisyntax{clojure}
(def paths {:a [:b :c]
            :b [:c]
            :c [:a]})

(vec (search/simple-paths paths [:a]))
;; => [{:state :a} {:state :b} {:state :c} {:state :c}]
\endtt

To see each route prefix explicitly, add the path extension:

\begtt \hisyntax{clojure}
(mapv :path
      (search/simple-paths paths [:a]
                           {:extensions [search/path-extension]}))
;; => [[:a] [:a :b] [:a :b :c] [:a :c]]
\endtt

The default path identity is \code{:state}; \code{:visit-key} can select a~
different identity (see section~\ref[identity]{Visit identity}), and
\code{:visit-key nil} disables the per-route repetition check. Simple-path
results are lazy; for example, a~bounded prefix of an infinite chain can be
consumed:

\begtt \hisyntax{clojure}
(mapv :state
      (take 4 (search/simple-paths (fn [n] [(inc n)]) [0])))
;; => [0 1 2 3]
\endtt

A~finite graph can still have many simple paths.

\secc Post-order traversal

\code{postorder} returns each state after its successors, in depth-first
post-order. When edges point from a~task to its prerequisites, prerequisites
therefore appear first. For example, this gives a~valid build order:

\begtt \hisyntax{clojure}
(vec (search/postorder {:package [:compile]
                          :compile [:parse]
                          :parse []}
                         [:package]))
;; => [{:state :parse} {:state :compile} {:state :package}]
\endtt

On a~cyclic graph, back edges are ignored for traversal and the search
terminates with DFS post-order when the reachable graph and its successor
collections are finite; that order is a~valid topological ordering only when
the reachable graph is acyclic.
Post-order traversal requires a~non-\code{nil} \code{:visit-key}.

Post-order traversal is lazy between completed visits, but it finishes
processing a~visit's successors before returning that visit. Taking the first
result does not require enumerating every successor. Here, the root has
infinitely many leaf successors, but the first completed visit is available
immediately:

\begtt \hisyntax{clojure}
(defn dependencies [task]
  (if (= task :root)
    (map #(vector :task %) (range))
    []))

(mapv :state (take 1 (search/postorder dependencies [:root])))
;; => [[:task 0]]
\endtt

The library decodes transitions and runs extension callbacks one at a~time.
A~chunked lazy successor collection returned by the space may itself compute
several raw successors when its first item is requested.

Post-order traversal also supports callbacks for processing completed
successors and pruning the remaining ones. These \code{:exit} and
\code{:more?} hooks are specific to post-order traversal; their examples
appear in section~\ref[transitions]{Additional visit and transition data}.

\secc State propagation

\code{propagate} is a~different kind of graph computation: instead of
discarding later arrivals at a~state, it joins their information and
processes that state again when the joined value changes. This is useful for
analyses that circulate information through cycles, such as reaching
definitions in a~compiler control-flow graph.

The space receives a~visit with \code{:state} and its current \code{:value};
transitions carry the destination and the contribution to join there.
\code{:join} describes how information combines, and \code{:bottom} supplies
the value for a~state with no previous contribution; it should be an
identity for the join. State map keys are identities, not visit keys. The
keys may be any Clojure value, including \code{nil} and \code{false}. The
result is a~lazy sequence of progress records. The first contains the initial
\code{:values} map and an empty \code{:changes} vector; it has no
\code{:processed} key because no state has been processed. Each later record
contains the updated map, the state just \code{:processed}, and the values
changed by its outgoing transitions. This distinction also makes the record
for processing a~\code{nil} state unambiguous. Use a~prefix to observe
intermediate progress or \code{last} to obtain the final map when propagation
terminates.

This small example shows the basic API: the value at \code{:start} passes
unchanged to its successor.

\begtt \hisyntax{clojure}
(require '[clojure.set :as set])

(def tiny-space {:start [:next]
                 :next []})

(:values (last (search/propagate tiny-space
                                 {:start #{:ready}}
                                 {:bottom #{}
                                  :join set/union})))
;; => {:start #{:ready}, :next #{:ready}}
\endtt

Unlike the search algorithms, propagation intentionally revisits a~state
when its value changes. For order-independent results, the join should be
associative, commutative, and idempotent, and transfer functions should be
monotone in the information order. Termination depends on the problem: the
finite set of facts in the worked loop analysis can only grow to a~fixed
point, but arbitrary spaces and value domains need not stabilize. These are
caller assumptions; the library does not validate them. Each expanded state
must return a~finite successor collection. Ordinary shortest-path costs do
not need this operation; \code{best-first} already handles them.

\secc Best-first search

\code{best-first} emits visits in increasing priority. By default its
\code{:distance} is the sum of transition weights, with each edge costing
one unless \code{:decode-edge} supplies a~\code{:weight}. It is therefore
Dijkstra's search when weights are finite and non-negative:

\begtt \hisyntax{clojure}
(def routes {:home [:slow :quick]
             :slow [:goal]
             :quick [:goal]
             :goal []})

(vec (search/best-first routes [:home]))
;; => [{:state :home :distance 0} {:state :slow :distance 1}
;;     {:state :quick :distance 1} {:state :goal :distance 2}]
\endtt

To account for different transition costs, decode each successor into a~
transition map with a~\code{:weight}:

\begtt \hisyntax{clojure}
(def weighted-routes
  {:home [[:slow 5] [:quick 1]]
   :quick [[:goal 2]]
   :slow [[:goal 1]]
   :goal []})

(mapv (juxt :state :distance)
      (search/best-first weighted-routes [:home]
                         {:decode-edge (fn [[state weight]]
                                         {:state state :weight weight})}))
;; => [[:home 0] [:quick 1] [:goal 3] [:slow 5]]
\endtt

An optional \code{:priority} function can order visits by something other than
distance, or add an estimate of the remaining cost as in A* search. A custom
priority does not automatically calculate \code{:distance}; include
\code{search/distance-extension} explicitly when the priority or results need
it. For example, a~route planner can combine the cost already travelled with
a~lower-bound estimate to the goal:

\begtt \hisyntax{clojure}
(def roads {:home [[:hill 1] [:river 2]]
            :hill [[:goal 5]]
            :river [[:goal 2]]
            :goal []})
(def estimate {:home 3 :hill 4 :river 2 :goal 0})

(->> (search/best-first roads [:home]
                        {:decode-edge (fn [[state weight]]
                                        {:state state :weight weight})
                         :extensions [search/path-extension
                                      search/distance-extension]
                         :priority (fn [{:keys [state distance]}]
                                     (+ distance (estimate state)))})
(filter #(= :goal (:state %)))
first
:path)
;; => [:home :river :goal]
\endtt

With the default \code{:visit-key} of \code{:state}, best-first settles each
state key at most once, using the best priority offered for it before
settlement. For A*, a~consistent heuristic with value zero at the goal is
sufficient for this behavior. An admissible but inconsistent heuristic may
need a~state to be explored again after a~better route is found. When the
heuristic depends only on the state, use a~\code{:visit-key} that distinguishes
different path costs, for example \code{(juxt :state :distance)}; this lets
best-first treat each state and cost pair as a~separate search position.
This may greatly increase the search and may prevent termination on cyclic
spaces. It is not a~general reopening mechanism for priorities that depend on
other visit context. Transition weights must be finite and non-negative.
Equal priorities are emitted in FIFO enqueue order; when starts or
successors come from an unordered collection, their tie order may vary.
Zero-weight edges are supported, but on an infinite space an infinite region
of equal or lower priority can indefinitely delay a~positive-cost result.
The usual A* completeness and optimality guarantees assume bounded branching,
an admissible heuristic, and edge costs bounded below by a~positive value;
with multiple-path pruning, consistency is also needed for the usual
optimality guarantee.
Best-first requires a~finite collection of starting states and finite
successor collections for each expanded visit. The default distance priority
automatically includes \code{distance-extension}; custom priorities do not.

Best-first is lazy as well. Since visits are emitted in priority order, a~
finite prefix can be consumed even when the search space is infinite:

\begtt \hisyntax{clojure}
(mapv :state
      (take 4 (search/best-first (fn [n] [(inc n)]) [0])))
;; => [0 1 2 3]
\endtt

\secc[identity] Visit identity

By default, DFS, BFS, post-order traversal, and best-first treat visits with
equal \code{:state} values as the same search position. This avoids exploring a~
state repeatedly and, for cyclic graphs, prevents traversal from following
cycles forever. More precisely, the visit key determines which arrivals an
algorithm treats as duplicates; later visits with a~key already seen are
skipped rather than combined with the earlier visit. The exact first arrival
depends on the search: DFS order for DFS, minimum edge-count depth for BFS in
a~static space, and minimum accumulated path cost for best-first with its
default distance priority and non-negative weights.

Sometimes context matters: the same state reached at different depths may
need to be explored twice. Choose a~\code{:visit-key} that includes that
context:

\begtt \hisyntax{clojure}
(def context-routes {:start [:near :detour]
                     :near [:goal]
                     :detour [:middle]
                     :middle [:goal]
                     :goal []})

(mapv (juxt :state :depth)
      (search/dfs context-routes [:start]
                  {:extensions [search/depth-extension]
                   :visit-key (juxt :state :depth)}))
;; => [[:start 0] [:near 1] [:goal 2]
;;     [:detour 1] [:middle 2] [:goal 3]]
\endtt

Sometimes context should affect which successors are allowed, while each
state should still be expanded only once. For example, explore links up to
two steps from the start, expanding each location at its shallowest arrival.
BFS reaches a~state at minimum depth, and the default \code{:visit-key}
discards later arrivals to that state:

\begtt \hisyntax{clojure}
(def depth-limited-routes
  {:start [:junction :detour]
   :detour [:junction]
   :junction [:goal]
   :goal [:finish]
   :finish []})

(defn within-depth [{:keys [state depth]}]
  (when (< depth 2)
    (depth-limited-routes state)))

(mapv (juxt :state :depth)
      (search/bfs within-depth [:start]
                  {:project identity
                   :extensions [search/depth-extension]}))
;; => [[:start 0] [:junction 1] [:detour 1] [:goal 2]]
\endtt

The second arrival at \code{:junction}, at depth two, is not returned or
expanded; the earlier arrival already expanded it with more steps remaining.
This is appropriate when the shallowest arrival is the one that should
govern exploration. If a~later arrival's context must also be explored, add
that context to \code{:visit-key}, as above.

The key function receives the full visit map. For DFS and BFS,
\code{:visit-key nil} disables global deduplication; on a~cycle, this can
produce an infinite result. \code{simple-paths} uses its key only to prevent
repeats within a~route; setting it to \code{nil} disables that check.
Identity affects correctness too: if equal keys can have different
successors, treating those visits as the same can change the search results.

For a~tree, disabling deduplication makes DFS visit each occurrence, like
\code{clojure.core/tree-seq}. The successor function supplies the children
directly, combining the roles of \code{tree-seq}'s \code{branch?} and
\code{children} functions; DFS still returns visit maps rather than bare
nodes. Here the two equal branch subtrees are distinct occurrences:

\begtt \hisyntax{clojure}
(def repeated-branch {:value :branch
                      :children [{:value :leaf}]})
(def tree {:value :root
           :children [repeated-branch repeated-branch]})
(defn values [nodes] (mapv :value nodes))

(values (map :state (search/dfs :children [tree])))
;; => [:root :branch :leaf]

(values (map :state (search/dfs :children [tree] {:visit-key nil})))
;; => [:root :branch :leaf :branch :leaf]

(values (tree-seq (comp seq :children) :children tree))
;; => [:root :branch :leaf :branch :leaf]
\endtt

With the default key, the second equal branch is skipped. Setting the key to
\code{nil} preserves both occurrences, matching \code{tree-seq}'s traversal.

\secc[transitions] Additional visit and transition data

Once an algorithm and its order are chosen, extensions can carry
problem-specific information in visits. \code{accumulate} builds an
extension from an initial-value function and a~step function. For example,
track remaining fuel along a~route:

\begtt \hisyntax{clojure}
(def fuel-extension
  (search/accumulate :fuel-left (constantly 10)
                     (fn [fuel transition]
                       (- fuel (:fuel-used transition)))))

(mapv (juxt :state :fuel-left)
      (search/dfs {:trailhead [[:ridge 3]]
                   :ridge []}
                  [:trailhead]
                  {:decode-edge (fn [[state fuel-used]]
                                  {:state state :fuel-used fuel-used})
                   :extensions [fuel-extension]}))
;; => [[:trailhead 10] [:ridge 7]]
\endtt

By default, a~space function receives only the current state. Set
\code{:project identity} to give it the full visit, including extension
values. This is useful when legal moves depend on search context:

\begtt \hisyntax{clojure}
(defn bounded-successors [{:keys [state depth]}]
  (when (< depth 2)
    [(inc state)]))

(vec (search/dfs bounded-successors [0]
                 {:project identity
                  :extensions [search/depth-extension]}))
;; => [{:state 0 :depth 0} {:state 1 :depth 1} {:state 2 :depth 2}]
\endtt

The \code{:expand?} option can prune a~visit without removing it from the
result. It receives the full visit; when it returns false, the space is not
called for that visit and no successors are generated:

\begtt \hisyntax{clojure}
(mapv (juxt :state :depth)
      (search/dfs {0 [1] 1 [2] 2 [3] 3 []}
                  [0]
                  {:extensions [search/depth-extension]
                   :expand? #(< (:depth %) 2)}))
;; => [[0 0] [1 1] [2 2]]
\endtt

Post-order traversal offers additional callbacks for work that depends on
completed successors. An \code{:exit} callback can compute a~value after a~
visit's successors have been processed. For example, count the tasks in each
dependency subtree:

\begtt \hisyntax{clojure}
(def subtree-size
  {:exit (fn [_visit successors]
           {:size (inc (reduce + (map :size successors)))})})

(mapv (juxt :state :size)
      (search/postorder {:app [:lib :config]
                           :lib [:base]
                           :base []
                           :config []}
                          [:app]
                          {:extensions [subtree-size]}))
;; => [[:base 1] [:lib 2] [:config 1] [:app 4]]
\endtt

An extension's \code{:more?} callback can stop consuming successors based
on those already processed. Here, a~cached build result prevents exploring
the fallback build:

\begtt \hisyntax{clojure}
(def cache-or-build
  {:more? (fn [{:keys [state]} previous]
            (or (not= state :root)
                (not-any? :ready? previous)))
   :exit (fn [{:keys [state]} successors]
           (if (= state :root)
             {:selected (or (some #(when (:ready? %) (:state %)) successors)
                            :build)}
             {:ready? (= state :cached)}))})

(mapv #(select-keys % [:state :ready? :selected])
      (search/postorder {:root [:cached :build]
                           :cached []
                           :build []}
                          [:root]
                          {:extensions [cache-or-build]}))
;; => [{:state :cached, :ready? true}
;;     {:state :root, :selected :cached}]
\endtt

A~\code{:next} callback can use completed siblings while extending the next
visit. Here it records each dependency's position in its parent's successor
order:

\begtt \hisyntax{clojure}
(def sibling-position
  {:next (fn [parent _transition _visit]
           {:position (count (:previous parent))})})

(mapv (juxt :state :position)
      (take 2
            (search/postorder {:app [:lib :config]
                                 :lib []
                                 :config []}
                                [:app]
                                {:extensions [sibling-position]})))
;; => [[:lib 0] [:config 1]]
\endtt

The \code{:next} callback receives the parent visit with its completed
successors in \code{:previous}. The \code{:exit} callback receives those
successor visits; each carries \code{:via}, the transition from this parent,
and \code{:back-edge}, true if it refers to a~visit still being explored.
Back edges have no exit values. This lets an exit callback inspect cycles
without following them again:

\begtt \hisyntax{clojure}
(def edge-report
  {:exit (fn [_visit successors]
           {:edges (mapv (juxt :state (comp :kind :via) :back-edge)
                         successors)})})

(mapv (juxt :state :edges)
      (search/postorder {:a [[:b :depends-on]]
                           :b [[:a :depends-on]]}
                          [:a]
                          {:decode-edge (fn [[state kind]]
                                          {:state state :kind kind})
                           :extensions [edge-report]}))
;; => [[:b [[:a :depends-on true]]]
;;     [:a [[:b :depends-on false]]]]
\endtt

An edge can carry information in addition to its destination, such as a~
label or weight. Use \code{:decode-edge} to turn each raw successor into a~
transition map containing \code{:state}; extensions can inspect the other
fields. The built-in \code{trail-extension} records decoded transitions
along a~route:

\begtt \hisyntax{clojure}
(def weighted-graph
  {:start [[:next 7]]
   :next []})

(mapv #(select-keys % [:state :trail])
      (search/dfs weighted-graph [:start]
                  {:decode-edge (fn [[state weight]]
                                  {:state state :weight weight})
                   :extensions [search/trail-extension]}))
;; => [{:state :start :trail [{:state :start}]}
;;     {:state :next :trail [{:state :start}
;;                           {:state :next :weight 7}]}]
\endtt

Without a~decoder, each raw successor is treated as a~state, even if it is
a~map. The trail starts with \code{\{:state start\}}; later entries are the
decoded transitions. DFS, BFS, and \code{simple-paths} do not interpret edge
weights. Best-first's default distance priority uses \code{:weight} to calculate
\code{:distance}; custom priorities need \code{search/distance-extension} to
calculate distance;
post-order traversal exposes each incoming transition as \code{:via}.
For option defaults and termination behavior, see
section~\ref[options]{Options and termination}.

\secc[options] Options and termination

The following options are common to the graph searches. Best-first also
accepts \code{:priority}. Post-order traversal supports additional extension
callbacks, described above.

\begitems
* \code{:project} chooses what is passed to the space, applied to the full
  visit. It defaults to \code{:state}.
* \code{:expand?} decides whether to request successors for a~visit, applied
  to the full visit. It defaults to a~function that accepts every visit; a~
  rejected visit is still returned.
* \code{:decode-edge} optionally converts each raw successor to a~transition
  map containing \code{:state}. Without a~decoder, the raw successor is the
  state.
* \code{:extensions} is an ordered collection of maps, empty by default.
  \code{:start}, \code{:next}, and (for post-order traversal)
  \code{:exit} callbacks return maps of visit fields to add or replace.
  Post-order traversal's \code{:more?} callbacks are checked before
  consuming and decoding the next successor; all must allow it for traversal
  to proceed.
  This check can also happen once when a~visit has no successors.
* \code{:visit-key} selects the key used for duplicate detection, applied to
  the full visit. It defaults to \code{:state}; its meaning and restrictions
  vary by algorithm.
  Post-order traversal and best-first require a~non-\code{nil} key.
* Best-first also accepts \code{:priority}, a~function of the visit that
  determines queue order. It defaults to \code{:distance}, which automatically
  enables \code{distance-extension}. A custom priority does not add
  \code{:distance}; add \code{search/distance-extension} to \code{:extensions}
  if it needs cumulative distance. Priority values must be mutually comparable.
\enditems

\code{simple-paths} uses \code{:visit-key} to prevent repeats on each route,
rather than to deduplicate the entire search.

\code{propagate} uses \code{:project}, \code{:expand?}, and
\code{:decode-edge}; \code{:bottom} and \code{:join} are required. Its map
keys identify states, so it does not use \code{:visit-key} or
\code{:extensions}.

To exhaust a~result sequence after finitely many visits, the start collection
and the reachable search must be finite, and calls to the space and callbacks
must return.
Disabling deduplication on a~cycle can make DFS or BFS produce visits forever.
Simple-path enumeration can still be large even on a~finite graph because
there may be many distinct routes.

For an infinite space, DFS can follow an infinite branch forever, preventing
later siblings or starts from being reached. BFS reaches every visit at
finite depth when the start collection is finite and each visit has finitely
many successors; an infinite successor collection can prevent deeper levels
from being reached.
Post-order traversal must finish a~visit's successors before yielding it.
Best-first requires finite starts and finite successor collections for each
expanded visit. DFS, BFS, simple-path enumeration, and best-first results are
lazy, so use a~bounded operation such as \code{take} or \code{first} with
\code{filter} when only part of a~result is needed. Post-order traversal is
lazy between completed visits, but a~parent is not returned until its
successors have been processed.

The space must be callable on the projected value and return \code{nil} or a~
collection of successors. \code{:project}, \code{:visit-key}, and
\code{:expand?} must be callable; \code{:visit-key} may be \code{nil} only
where allowed above. When supplied, \code{:decode-edge} must return a~map
containing \code{:state}. Extension callbacks must return maps. Visit keys
should have stable equality and hashing. These input contracts are not
exhaustively validated at runtime.

\sec[examples] Worked problems

These examples form a~learning path. Each starts with a~problem, models its
states and transitions, then chooses a~library operation that fits the
question. They reuse ordinary Clojure functions and collections; the library
supplies the traversal or propagation behavior and carries context or
accumulated information as needed. Evaluate the examples in order in one REPL
session, after loading \code{state-space.core} as shown above.

For an exercise-led route, try the next two problems in order. The first
combines capabilities across a~service dependency graph; the second adds
transfer rules and a~cycle, so a~state must be processed again when its
information changes. The remaining problems offer other applications and
can be explored in any order.

For each problem, pause after reading the question and work out a~state
representation and an approach before continuing to the solution. The
solutions are written as runnable REPL examples.

Many examples use \code{first} after \code{filter} to stop as soon as a~
useful visit is found. Because the search results are lazy, this does not
need to finish exploring the rest of the space. Examples that collect every
answer, such as subset generation, intentionally consume the full result.

\secc[exercise-union] Find capabilities required by a~service

\noindent{\bf Problem.} A~storefront depends on a~catalog API and an account
API. The catalog API depends on a~database that provides \code{:sql}; the
account API depends on an identity provider that provides \code{:oauth}.
Which capabilities are required by each service, including the storefront?

\smallskip
\noindent{\bf Try it.} Draw edges from each dependency to the service that
uses it. Start with the capabilities provided by the database and identity
provider, then determine how to combine capabilities when dependencies meet.

\smallskip
\noindent{\bf Solution.} Use sets for capabilities and set union when
contributions meet. With the default projection and edge decoding, an
adjacency map describes how capabilities flow from dependencies to services:

\begtt \hisyntax{clojure}
(require '[clojure.set :as set])

(def service-dependencies
  {:catalog-db       [:catalog-api]
   :identity-provider [:account-api]
   :catalog-api      [:storefront]
   :account-api      [:storefront]
   :storefront       []})

(def values
  (:values
   (last (search/propagate service-dependencies
                           {:catalog-db #{:sql}
                            :identity-provider #{:oauth}}
                           {:bottom #{}
                            :join set/union}))))
;; => {:catalog-db #{:sql}
;;     :identity-provider #{:oauth}
;;     :catalog-api #{:sql}
;;     :account-api #{:oauth}
;;     :storefront #{:sql :oauth}}
\endtt

The two branches bring different requirements to the storefront, where set
union combines them. Each service's value is the set of capabilities it
inherits from its dependencies. The data is deliberately small; the same
model can be used to inspect larger dependency graphs.

\secc[exercise-reaching] Analyze definitions around a~loop

\noindent{\bf Problem.} A~program has an entry block, a~loop header, a~loop
body, and an exit. The entry assigns to \code{x}; the body assigns to
\code{x} again before returning to the header. At each block, determine
which assignments to \code{x} may reach it. A~definition is a~pair of a~
variable and its assignment label.

\smallskip
\noindent{\bf Try it.} Start with no definitions at the entry. At each
block, remove definitions it overwrites and add definitions it generates.
When paths meet, combine the possible definitions. What changes when the
loop's definition travels back to the header?

\smallskip
\noindent{\bf Solution.} Represent each block by its successors and its
\code{gen}/\code{kill} sets. The transfer function applies those sets before
sending a~contribution along every outgoing edge:

\begtt \hisyntax{clojure}
(def blocks
  {:entry  {:successors [:header] :gen #{[:x :initial]} :kill #{}}
   :header {:successors [:body :exit] :gen #{} :kill #{}}
   :body   {:successors [:header]
            :gen #{[:x :loop]}
            :kill #{[:x :initial] [:x :loop]}}
   :exit   {:successors [] :gen #{} :kill #{}}})

(defn transfer [{:keys [state value]}]
  (let [{:keys [successors gen kill]} (blocks state)
        outgoing (into gen (remove kill value))]
    (map (fn [target] {:state target :value outgoing}) successors)))

(def progress
  (search/propagate transfer
                    {:entry #{}}
                    {:project identity
                     :decode-edge identity
                     :bottom #{}
                     :join set/union}))

(:values (last progress))
;; => {:entry #{}
;;     :header #{[:x :initial] [:x :loop]}
;;     :body #{[:x :initial] [:x :loop]}
;;     :exit #{[:x :initial] [:x :loop]}}
\endtt

The body kills the entry definition and generates the loop definition on
each pass. At the header, contributions from the entry path and the back
edge are joined, so the header's value changes and it is processed again.
The finite set of possible definitions grows only to a~fixed point. For
arbitrary value domains, propagation need not terminate; the join and
transfer behavior must satisfy the assumptions described above.

\secc Which route does depth-first search find first?

\noindent{\bf Problem.} A~building has rooms \code{:lobby},
\code{:stairs}, \code{:landing}, \code{:gallery}, and \code{:exit}. Its
directed passages, in the order inspected from each room, are
\code{:lobby -> :stairs}, \code{:lobby -> :gallery},
\code{:stairs -> :landing}, \code{:landing -> :exit}, and
\code{:gallery -> :exit}. Starting at \code{:lobby}, what route to
\code{:exit} does depth-first search find first?

\smallskip
\noindent{\bf Solution.} Represent each room as a~state and each passage as
a~successor. Record the route and stop at the first visit to the exit:

\begtt \hisyntax{clojure}
(def building
  {:lobby [:stairs :gallery]
   :stairs [:landing]
   :landing [:exit]
   :gallery [:exit]
   :exit []})

(first (filter #(= :exit (:state %))
               (search/dfs building [:lobby]
                           {:extensions [search/path-extension]})))
;; => {:state :exit, :path [:lobby :stairs :landing :exit]}
\endtt

DFS takes the stairs branch first and returns a~three-passage route, even
though the gallery gives a~shorter route. DFS answers which route it finds
first, not which route is shortest; a~later example uses BFS for shortest
routes.

\secc Generate every subset

\noindent{\bf Problem.} List every subset of the four distinct items
\code{:a}, \code{:b}, \code{:c}, and \code{:d}, including the empty set and
the set containing all four. Represent each subset as a~vector in the
items' original order.

\smallskip
\noindent{\bf Solution.} At each step, decide whether to leave the next
item out or add it to the partial result. DFS explores this decision tree;
the index in the state tells us when all decisions have been made:

\begtt \hisyntax{clojure}
(def items [:a :b :c :d])

(defn subset-decisions [[i chosen]]
  (when (< i (count items))
    [[(inc i) chosen]
     [(inc i) (conj chosen (items i))]]))

(->> (search/dfs subset-decisions [[0 []]])
     (map :state)
     (filter #(= (count items) (first %)))
     (map second)
     vec)
;; => [[] [:d] [:c] [:c :d]
;;     [:b] [:b :d] [:b :c] [:b :c :d]
;;     [:a] [:a :d] [:a :c] [:a :c :d]
;;     [:a :b] [:a :b :d] [:a :b :c] [:a :b :c :d]]
\endtt

The search state includes both the partial answer and how many choices have
been considered. This lets DFS distinguish unfinished decisions from
completed subsets.

\secc Place eight queens without conflicts

\noindent{\bf Problem.} Place eight queens on a~chessboard so that no two
share a~row, column, or diagonal. How many solutions are there?

\smallskip
\noindent{\bf Solution.} Place one queen per row. A~partial placement is a~
vector of column numbers, so row conflicts are impossible by construction.
Reject a~column as soon as it conflicts with a~queen already placed; DFS
then explores only safe partial boards:

\begtt \hisyntax{clojure}
(defn safe-column? [placed column]
  (let [row (count placed)]
    (every? (fn [[previous-row previous-column]]
              (and (not= column previous-column)
                   (not= (- row previous-row)
                         (Math/abs (- column previous-column)))))
            (map-indexed vector placed))))

(defn queen-placements [placed]
  (when (< (count placed) 8)
    (for [column (range 8)
          :when (safe-column? placed column)]
      (conj placed column))))

(count (filter #(= 8 (count (:state %)))
               (search/dfs queen-placements [[]])))
;; => 92
\endtt

The partial column vector is the search state. Pruning unsafe choices as
soon as they appear avoids exploring complete boards that could never be
solutions.

\secc Find a~way around a~blocked square

\noindent{\bf Problem.} A~robot is on a $5 \times 5$ grid whose positions
are \code{[x y]} for \code{x} and \code{y} from 0 through 4. It starts at
\code{[0 2]} and can move one square horizontally or vertically. It cannot
leave the grid or enter any of the blocked squares \code{[2 0]},
\code{[2 1]}, \code{[2 2]}, and \code{[2 3]}. What is the fewest moves
needed to reach \code{[4 2]}?

\smallskip
\noindent{\bf Solution.} A~state is the robot's position. Generate adjacent
positions and keep those on the grid and outside the blocked set. BFS
explores states in order of depth, so the depth extension reports the
fewest moves:

\begtt \hisyntax{clojure}
(def blocked-squares #{[2 0] [2 1] [2 2] [2 3]})

(defn grid-moves [[x y]]
  (filter (fn [[nx ny]]
            (and (<= 0 nx 4)
                 (<= 0 ny 4)
                 (not (contains? blocked-squares [nx ny]))))
          [[(inc x) y] [(dec x) y] [x (inc y)] [x (dec y)]]))

(->> (search/bfs grid-moves [[0 2]]
                 {:extensions [search/depth-extension]})
     (filter #(= [4 2] (:state %)))
     first
     :depth)
;; => 8
\endtt

This is an implicit graph: the successor function generates legal moves
when a~position is explored instead of storing the whole graph.

\secc Find each cell's distance to the nearest zero

\noindent{\bf Problem.} A~grid contains zeros and ones. For every cell,
find its Manhattan distance to the nearest zero. For example, in
\code{["000" "010" "111"]}, the middle cell in the second row is one step
from a~zero, while the bottom-right cell is two steps from one.

\smallskip
\noindent{\bf Solution.} Start one BFS from every zero at once. BFS visits
each cell at its minimum distance from any start; collect those distances
into a~map and arrange them back into rows:

\begtt \hisyntax{clojure}
(def grid ["000" "010" "111"])

(def starts
  (for [row (range (count grid))
        col (range (count (first grid)))
        :when (= \0 (get-in grid [row col]))]
    [row col]))

(defn neighbours [[row col]]
  (filter #(get-in grid %)
          [[(dec row) col] [(inc row) col]
           [row (dec col)] [row (inc col)]]))

(def distance
  (into {} (map (juxt :state :depth))
        (search/bfs neighbours starts
                    {:extensions [search/depth-extension]})))

(mapv (fn [row]
        (mapv #(distance [row %]) (range (count (first grid)))))
      (range (count grid)))
;; => [[0 0 0] [0 1 0] [1 2 1]]
\endtt

The multiple starting states make each cell's first BFS arrival its distance
to the nearest zero.

\secc Check whether a~graph is bipartite

\noindent{\bf Problem.} An undirected graph has two disconnected components:
a~four-cycle and a~triangle. Can its vertices be split into two groups so
that every edge connects different groups?

\smallskip
\noindent{\bf Solution.} In each component, BFS depth parity gives a~two-coloring
candidate. The graph is bipartite exactly when every edge joins vertices of
opposite parity. Start a~new BFS at every still-uncoloured component so a~
disconnected triangle cannot be missed:

\begtt \hisyntax{clojure}
(def graph
  {:a [:b :d] :b [:a :c] :c [:b :d] :d [:a :c]
   :x [:y :z] :y [:x :z] :z [:x :y]})

(def nodes (into (set (keys graph)) (mapcat identity (vals graph))))

(def adjacent
  (reduce (fn [g [from to]]
            (-> g (update from conj to) (update to conj from)))
          (zipmap nodes (repeat []))
          (for [[from tos] graph, to tos] [from to])))

(loop [uncoloured nodes]
  (if-let [root (first uncoloured)]
    (let [component
          (into {} (map (fn [{:keys [state depth]}]
                          [state (mod depth 2)]))
                (search/bfs adjacent [root]
                            {:extensions [search/depth-extension]}))]
      (if (some (fn [[node colour]]
                  (some #(= colour (component %)) (adjacent node)))
                component)
        false
        (recur (reduce disj uncoloured (keys component))))
    true))
;; => false
\endtt

The square is bipartite, but the separate triangle contains an edge whose
endpoints have the same parity. Checking every component is essential.

\secc Browse links up to two clicks away

\noindent{\bf Problem.} A~website has pages \code{:home}, \code{:about},
\code{:shop}, \code{:team}, \code{:item}, and \code{:reviews}. Its only
links, in the order inspected, are \code{:home -> :about},
\code{:home -> :shop}, \code{:about -> :team}, \code{:about -> :shop},
\code{:shop -> :about}, \code{:shop -> :team}, \code{:shop -> :item},
\code{:team -> :home}, \code{:team -> :reviews}, and
\code{:item -> :reviews}. Starting at \code{:home}, list every page
reachable by following at most two links, with its minimum number of clicks.
Include \code{:home} itself, but do not follow links from pages reached in
two clicks.

\smallskip
\noindent{\bf Solution.} Use BFS and add the depth extension. Here the space
function reads both the current page and the search depth: it follows links
only when fewer than two clicks have been made. Because BFS reaches each page
at its minimum depth and the default visit identity is just \code{:state},
the first arrival is the one with the most clicks still available:

\begtt \hisyntax{clojure}
(def website
  {:home [:about :shop]
   :about [:team :shop]
   :shop [:about :team :item]
   :team [:home :reviews]
   :item [:reviews]
   :reviews []})

(defn website-links [{:keys [state depth]}]
  (when (< depth 2)
    (website state)))

(mapv (juxt :state :depth)
      (search/bfs website-links [:home]
                  {:project identity
                   :extensions [search/depth-extension]}))
;; => [[:home 0] [:about 1] [:shop 1] [:team 2] [:item 2]]
\endtt

The space receives the full visit because of \code{:project identity}, so
successor generation depends on \code{:depth} as well as \code{:state}.
Visits are still deduplicated by \code{:state}: BFS guarantees that the first
time a~page is expanded is its shallowest arrival. Pages at the two-click
limit are included, but produce no successors; cycles and repeated links do
not produce duplicate visits.

For this simple cutoff, \code{:expand?} is more direct: keep the website as a~
state-based space and let the option prevent expansion at the limit. Both
versions return the same visits. Use a~search-aware space when context
determines which successors are available, rather than whether to request any
successors at all.

\begtt \hisyntax{clojure}
(mapv (juxt :state :depth)
      (search/bfs website [:home]
                  {:extensions [search/depth-extension]
                   :expand? #(< (:depth %) 2)}))
;; => [[:home 0] [:about 1] [:shop 1] [:team 2] [:item 2]]
\endtt

\secc Fewest moves across a~chessboard

\noindent{\bf Problem.} A~knight starts at \code{[0 0]} on a~standard
$8 \times 8$ chessboard whose squares are \code{[x y]} with coordinates
from 0 through 7. A~knight moves two squares in one direction and one square
perpendicular to that direction, staying on the board. What is the fewest
moves needed to reach the opposite corner, \code{[7 7]}?

\smallskip
\noindent{\bf Solution.} Treat each square as a~state. Generate the legal
knight moves from a~square, then use BFS because every move has equal cost.

\begtt \hisyntax{clojure}
(defn knight-moves [[x y]]
  (for [[dx dy] [[1 2] [2 1] [-1 2] [-2 1]
                 [1 -2] [2 -1] [-1 -2] [-2 -1]]
        :let [next [(+ x dx) (+ y dy)]]
        :when (every? #(<= 0 % 7) next)]
    next))

(->> (search/bfs knight-moves [[0 0]]
                 {:extensions [search/depth-extension]})
     (filter #(= [7 7] (:state %)))
     first
     :depth)
;; => 6
\endtt

The successor function generates only legal squares; no graph needs to be
stored.

\secc Measure four litres with two jugs

\noindent{\bf Problem.} You have an unmarked three-litre jug and an unmarked
five-litre jug, both initially empty. You may fill either jug completely
from a~tap, empty either jug, or pour from one jug to the other until the
source is empty or the destination is full. What is the fewest actions
needed to leave exactly four litres in the five-litre jug?

\smallskip
\noindent{\bf Solution.} Represent a~state as the amount of water in both
jugs. Each allowed action creates a~successor state. BFS finds a~solution
with the fewest actions, and the path extension shows the states along it:

\begtt \hisyntax{clojure}
(defn jug-moves [[a b]]
  (let [capacity-a 3 capacity-b 5
        pour-ab (min a (- capacity-b b))
        pour-ba (min b (- capacity-a a))]
    [[capacity-a b] [a capacity-b] [0 b] [a 0]
     [(- a pour-ab) (+ b pour-ab)]
     [(+ a pour-ba) (- b pour-ba)]]))

(->> (search/bfs jug-moves [[0 0]]
                 {:extensions [search/path-extension]})
     (filter (fn [{[_ b] :state}] (= 4 b)))
     first
     :path)
;; => [[0 0] [0 5] [3 2] [0 2] [2 0] [2 5] [3 4]]
\endtt

The full state matters: the same amount in one jug can lead to different
future moves depending on the other jug's contents.

\secc Help the missionaries and cannibals cross

\noindent{\bf Problem.} Three missionaries and three cannibals are on one
bank of a~river. A~boat carries one or two people. Whenever missionaries
are present on a~bank, they must not be outnumbered by cannibals. What is
the fewest crossings that take everyone to the opposite bank?

\smallskip
\noindent{\bf Solution.} Record how many missionaries and cannibals remain
on the starting bank, plus which bank has the boat. A~crossing is legal only
when both banks remain safe. BFS finds the minimum number of crossings:

\begtt \hisyntax{clojure}
(defn safe-bank? [missionaries cannibals]
  (or (zero? missionaries)
      (>= missionaries cannibals)))

(defn safe-state? [[left-m left-c _boat]]
  (and (<= 0 left-m 3)
       (<= 0 left-c 3)
       (safe-bank? left-m left-c)
       (safe-bank? (- 3 left-m) (- 3 left-c))))

(defn crossings [[missionaries cannibals boat]]
  (let [direction (if (zero? boat) -1 1)]
    (for [[move-m move-c] [[1 0] [2 0] [0 1] [0 2] [1 1]]
          :let [next [(+ missionaries (* direction move-m))
                      (+ cannibals (* direction move-c))
                      (- 1 boat)]]
          :when (safe-state? next)]
      next)))

(->> (search/bfs crossings [[3 3 0]]
                 {:extensions [search/depth-extension]})
     (filter #(= [0 0 1] (:state %)))
     first
     :depth)
;; => 11
\endtt

Checking both banks after every move keeps unsafe states out of the search.
The boat side is part of the state because it determines the direction of
the next crossing.

\secc Make a~number from the fewest perfect squares

\noindent{\bf Problem.} Starting from 12, repeatedly subtract a~positive
perfect square no larger than the current number. What is the fewest
subtractions needed to reach zero?

\smallskip
\noindent{\bf Solution.} Use the remaining number as the state, with one
successor for each legal square subtraction. BFS finds a~route with the
fewest subtractions:

\begtt \hisyntax{clojure}
(defn square-subtractions [remaining]
  (for [root (range 1 (inc (long (Math/sqrt remaining))))]
    (- remaining (* root root))))

(->> (search/bfs square-subtractions [12]
                 {:extensions [search/path-extension]})
     (filter #(zero? (:state %)))
     first
     :path)
;; => [12 8 4 0]
\endtt

Each step subtracts a~square; the path therefore represents
\code{12 = 4 + 4 + 4}.

\secc Open a~combination lock in the fewest turns

\noindent{\bf Problem.} A~combination lock has four decimal wheels,
initially set to \code{"0000"}. One turn moves one wheel up or down by one,
wrapping between 0 and 9. The combinations \code{"0201"}, \code{"0101"},
\code{"0102"}, \code{"1212"}, and \code{"2002"} are blocked and cannot be
entered. What is the fewest turns needed to reach \code{"0202"}?

\smallskip
\noindent{\bf Solution.} A~four-character string is a~state. Generate the
eight possible single-wheel turns, wrapping each digit with modular
arithmetic, and omit blocked combinations. BFS finds the minimum number of
turns:

\begtt \hisyntax{clojure}
(def blocked-codes #{"0201" "0101" "0102" "1212" "2002"})

(defn lock-turns [code]
  (for [wheel (range 4)
        direction [-1 1]
        :let [digit (Character/digit ^char (nth code wheel) 10)
              next-digit (mod (+ digit direction) 10)
              next-code (str (subs code 0 wheel)
                             next-digit
                             (subs code (inc wheel)))]
        :when (not (blocked-codes next-code))]
    next-code))

(->> (search/bfs lock-turns ["0000"]
                 {:extensions [search/depth-extension]})
     (filter #(= "0202" (:state %)))
     first
     :depth)
;; => 6
\endtt

The search need not construct all 10,000 combinations in advance; it
generates legal turns only for combinations it explores.

\secc Solve a~sliding-tile puzzle

\noindent{\bf Problem.} A~two-row, three-column board contains tiles
\code{1} through \code{5} and one blank, represented by \code{0}. Represent
the board as a~vector in row-major order. The blank may swap places with a~
tile directly above, below, left, or right of it, provided that position is
on the board. Starting from \code{[4 1 3 0 2 5]}, find the fewest swaps
that produce \code{[1 2 3 4 5 0]}.

\smallskip
\noindent{\bf Solution.} Represent the board as a~vector in row-major
order. A~state is the whole board; generate successors by swapping the
blank with each adjacent position. BFS finds a~shortest solution:

\begtt \hisyntax{clojure}
(defn slide-blank [tiles]
  (let [blank (.indexOf ^java.util.List tiles 0)
        row (quot blank 3)
        col (rem blank 3)]
    (for [[dr dc] [[-1 0] [1 0] [0 -1] [0 1]]
          :let [r (+ row dr)
                c (+ col dc)]
          :when (and (<= 0 r 1)
                     (<= 0 c 2))]
      (assoc tiles blank
             (tiles (+ (* r 3) c))
             (+ (* r 3) c) 0))))

(->> (search/bfs slide-blank [[4 1 3 0 2 5]]
                 {:extensions [search/path-extension]})
     (filter #(= [1 2 3 4 5 0] (:state %)))
     first
     :path)
;; => [[4 1 3 0 2 5]
;;     [0 1 3 4 2 5]
;;     [1 0 3 4 2 5]
;;     [1 2 3 4 0 5]
;;     [1 2 3 4 5 0]]
\endtt

The path contains four swaps. Unlike the jug puzzle, here each move changes
the arrangement of the whole state rather than quantities in separate
containers.

\secc Push a~box to its goal

\noindent{\bf Problem.} In the small warehouse below, move the box marked
\code{\$} onto the goal marked \code{.}. The player may walk through empty
cells and push the box, but cannot move through walls or pull the box. What
is the fewest player moves?

\begtt
#######
#@    #
# $ . #
#     #
#######
\endtt

\smallskip
\noindent{\bf Solution.} A~state contains both the player's position and
the box's position. When the player steps into the box, that same move
pushes it one cell farther; reject the move if the box would hit a~wall.
BFS searches these joint configurations by number of player moves:

\begtt \hisyntax{clojure}
(def warehouse
  ["#######"
   "#@    #"
   "# $ . #"
   "#     #"
   "#######"])

(defn wall? [[row column]]
  (= \# (get-in warehouse [row column])))

(defn sokoban-moves [[player box]]
  (for [direction [[1 0] [-1 0] [0 1] [0 -1]]
        :let [next-player (mapv + player direction)]
        :when (not (wall? next-player))
        :let [next-box (if (= next-player box)
                         (mapv + box direction)
                         box)]
        :when (not (wall? next-box))]
    [next-player next-box]))

(->> (search/bfs sokoban-moves [[[1 1] [2 2]]]
                 {:extensions [search/depth-extension]})
     (filter #(= [2 4] (second (:state %))))
     first
     :depth)
;; => 3
\endtt

Combining the two positions in one state lets the ordinary visited-state
check distinguish configurations with the box in different places.

\secc Plan around the Sussman anomaly

\noindent{\bf Problem.} Three blocks are arranged in two stacks, with the
bottom block first: \code{[:a :c]} and \code{[:b]}. Move one clear block at
a~time, either onto the table or onto another clear block, until the stack
\code{[:c :b :a]} is built. What is the fewest moves?

\smallskip
\noindent{\bf Solution.} Represent a~world as a~set of stacks. A~move
removes a~stack's top block and places it on the table or on another stack's
top. BFS searches complete world states, so it can find a~plan that first
undoes part of the initial arrangement:

\begtt \hisyntax{clojure}
(defn block-moves [stacks]
  (for [from stacks
        :let [block (peek from)
              without (cond-> (disj stacks from)
                        (seq (pop from)) (conj (pop from)))]
        target (cons nil (disj stacks from))]
    (if target
      (-> without (disj target) (conj (conj target block)))
      (conj without [block]))))

(let [plan (first (filter #(= #{[:c :b :a]} (:state %))
                          (search/bfs block-moves [#{[:a :c] [:b]}]
                                      {:extensions [search/depth-extension
                                                    search/path-extension]})))]
  {:depth (:depth plan)
   :path (mapv #(vec (sort-by first %)) (:path plan))})
;; => {:depth 3, :path [[[:a :c] [:b]] [[:a] [:b] [:c]]
;;                      [[:a] [:c :b]] [[:c :b :a]]]}
\endtt

The planner must temporarily move \code{:c} off \code{:a} so that it can
build \code{:c} on \code{:b}, then put \code{:a} on top. Treating the goals
“put \code{:c} on the table” and “put \code{:a} on \code{:b}” as independent
subtasks would miss that interaction.

\secc Turn a~word into another word

\noindent{\bf Problem.} The allowed dictionary consists of \code{cold},
\code{cord}, \code{card}, \code{ward}, \code{warm}, \code{worm},
\code{word}, and \code{corm}. Starting from \code{cold}, reach \code{warm}
by changing exactly one letter at a~time. Every word in the route,
including its start and end, must be in the dictionary. Find a~route with
the fewest changes.

\smallskip
\noindent{\bf Solution.} Regard each word as a~state and compute its
neighbours from the dictionary. BFS finds the shortest route:

\begtt \hisyntax{clojure}
(def words #{"cold" "cord" "card" "ward" "warm"
             "worm" "word" "corm"})

(defn one-letter-apart? [a b]
  (= 1 (count (filter true? (map not= a b)))))

(defn ladder-neighbours [word]
  (filter #(one-letter-apart? word %) words))

(->> (search/bfs ladder-neighbours ["cold"]
                 {:extensions [search/path-extension]})
     (filter #(= "warm" (:state %)))
     first
     :path)
;; => ["cold" "cord" "card" "ward" "warm"]
\endtt

The space is implicit, and lazy \code{filter} stops consuming BFS as soon as
it reaches the goal.

\secc List all routes without repeating a~place

\noindent{\bf Problem.} A~directed map has places \code{:home},
\code{:park}, \code{:market}, and \code{:museum}. Its only roads are
\code{:home -> :park}, \code{:home -> :market},
\code{:park -> :market}, \code{:market -> :park},
\code{:park -> :museum}, and \code{:market -> :museum}. List every route
from \code{:home} to \code{:museum} that visits no place more than once.

\smallskip
\noindent{\bf Solution.} Use \code{simple-paths}: unlike DFS, it keeps
separate routes that reach the same place, while preventing a~route from
looping back to a~place it has already visited. The path extension records
each route prefix:

\begtt \hisyntax{clojure}
(def city-map
  {:home [:park :market]
   :park [:market :museum]
   :market [:park :museum]
   :museum []})

(->> (search/simple-paths city-map [:home]
                          {:extensions [search/path-extension]})
     (filter #(= :museum (:state %)))
     (map :path)
     vec)
;; => [[:home :park :market :museum]
;;     [:home :park :museum]
;;     [:home :market :park :museum]
;;     [:home :market :museum]]
\endtt

This enumeration is useful when all routes matter, rather than only the
first route or a~shortest route. Even in a~finite graph, the number of
simple paths can grow quickly.

\secc Find all walks of exactly three steps

\noindent{\bf Problem.} A~directed graph has states \code{:a}, \code{:b},
and \code{:c}, with only the edges \code{:a -> :b}, \code{:b -> :a},
\code{:b -> :c}, and \code{:c -> :a}. List every walk of exactly three
steps starting at \code{:a}. A~walk may visit the same state more than once.

\smallskip
\noindent{\bf Solution.} Stop expanding visits at depth three, but keep
states distinct by route: disable the default global \code{:visit-key}
deduplication and record each path. Then select the visits at the required
depth:

\begtt \hisyntax{clojure}
(def walk-graph
  {:a [:b]
   :b [:a :c]
   :c [:a]})

(->> (search/bfs walk-graph [:a]
                 {:visit-key nil
                  :expand? #(< (:depth %) 3)
                  :extensions [search/depth-extension
                               search/path-extension]})
     (filter #(= 3 (:depth %)))
     (map :path)
     set)
;; => #{[:a :b :a :b] [:a :b :c :a]}
\endtt

This differs from \code{simple-paths}: here revisiting a~state is allowed,
because the question is about walks of a~fixed length, not cycle-free routes.

\secc Visit every node in the fewest steps

\noindent{\bf Problem.} A~robot moves along the edges of this undirected
star-shaped graph. It may start at any node and revisit nodes. What is the
fewest moves needed to visit all four nodes?

\smallskip
\noindent{\bf Solution.} The current node alone is not enough to describe
the search: future success also depends on which nodes have already been
visited. Use the pair \code{[node visited]} as the state and start a~BFS
from each possible node:

\begtt \hisyntax{clojure}
(def graph [[1 2 3]
            [0]
            [0]
            [0]])

(def all-nodes (set (range (count graph))))

(defn visit-next [[node visited]]
  (for [next (graph node)]
    [next (conj visited next)]))

(def starts
  (map (fn [node] [node #{node}])
       (range (count graph))))

(->> (search/bfs visit-next starts
                 {:extensions [search/depth-extension]})
     (filter #(= all-nodes (second (:state %))))
     first
     :depth)
;; => 4
\endtt

Because the visited set is part of state identity, reaching the same graph
node with different progress is not mistaken for a~duplicate visit.

\secc Check mutual exclusion across process interleavings

\noindent{\bf Problem.} Two processes use Peterson's mutual-exclusion
protocol. Can both enter their critical sections at once? What shortest
counterexample appears if the waiting condition is removed?

\smallskip
\noindent{\bf Solution.} A~state records each process's control location,
its request flag, and whose turn it is. At each step either process may
advance by one instruction. BFS can check that the correct protocol never
reaches a~state where both program counters are \code{3}, then return a~
shortest trace when the waiting condition is broken:

\begtt \hisyntax{clojure}
(def initial-system {:pc [0 0] :flag [false false] :turn 0})

(defn process-steps [wait?]
  (fn [{:keys [turn] :as system}]
    (for [i [0 1]
          :let [j (- 1 i)]]
      (case (get-in system [:pc i])
        0 (-> system (assoc-in [:flag i] true) (assoc-in [:pc i] 1))
        1 (-> system (assoc :turn j) (assoc-in [:pc i] 2))
        2 (if (and wait? (get-in system [:flag j]) (= turn j))
            system
            (assoc-in system [:pc i] 3))
        3 (assoc-in system [:pc i] 4)
        4 (-> system (assoc-in [:flag i] false)
              (assoc-in [:pc i] 0)))))

(defn violation [wait?]
  (first (filter #(= [3 3] (get-in % [:state :pc]))
                 (search/bfs (process-steps wait?) [initial-system]
                             {:extensions [search/depth-extension
                                           search/path-extension]}))))

(let [counterexample (violation false)]
  [(nil? (violation true))
   (:depth counterexample)
   (mapv #(get-in % [:pc]) (:path counterexample))])
;; => [true 6 [[0 0] [1 0] [2 0] [3 0] [3 1] [3 2] [3 3]]]
\endtt

The correct protocol preserves mutual exclusion for every interleaving.
Without waiting, the trace lets each process pass its check before either
enters, violating the safety property after six instructions.

\secc Choose the least expensive route

\noindent{\bf Problem.} A~travel network has directed trips with these
times in hours: from \code{:home} to \code{:museum} (11), \code{:village}
(4), or \code{:park} (1); from \code{:village} to \code{:museum} (7) or
\code{:hill} (2); from \code{:park} to \code{:museum} (8) or
\code{:village} (1); and from \code{:hill} to \code{:museum} (2). There
are no other trips. Which route from \code{:home} to \code{:museum} takes
the least time, and what is its total cost?

\smallskip
\noindent{\bf Solution.} Each edge stores its travel time. Decode an edge
into a~transition containing both its destination and weight; best-first
uses the accumulated distance to return the least-cost route:

\begtt \hisyntax{clojure}
(def travel-times
  {:home [[:museum 11] [:village 4] [:park 1]]
   :village [[:museum 7] [:hill 2]]
   :park [[:museum 8] [:village 1]]
   :hill [[:museum 2]]
   :museum []})

(->> (search/best-first travel-times [:home]
                        {:decode-edge (fn [[state weight]]
                                        {:state state :weight weight})
                         :extensions [search/path-extension]})
     (filter #(= :museum (:state %)))
     first
     (select-keys [:path :distance]))
;; => {:path [:home :park :village :hill :museum], :distance 6}
\endtt

Unlike BFS, best-first minimizes the sum of edge costs, not the number of
edges. Its default distance priority is Dijkstra's algorithm when edge
weights are non-negative.

\secc Find a~least-cost route with a~negative edge

\noindent{\bf Problem.} A~directed acyclic network has trips
\code{:s -> :a} (4), \code{:s -> :b} (2), \code{:a -> :t} (-3),
\code{:b -> :a} (-1), and \code{:b -> :t} (5). What is the least total
cost from \code{:s} to \code{:t}?

\smallskip
\noindent{\bf Solution.} Dijkstra's priority rule relies on non-negative
edge weights. In a~DAG, post-order traversal can instead evaluate the
cheapest remaining route from each node after its successors are known. The
decoded incoming edge supplies its weight:

\begtt \hisyntax{clojure}
(def dag
  {:s [{:state :a :weight 4} {:state :b :weight 2}]
   :a [{:state :t :weight -3}]
   :b [{:state :a :weight -1} {:state :t :weight 5}]
   :t []})

(def route-cost
  {:exit (fn [{:keys [state]} successors]
           {:cost (if (= state :t)
                    0
                    (apply min
                           (map #(+ (:weight (:via %)) (:cost %))
                                successors)))})})

(:cost (last (search/postorder dag [:s]
                                 {:decode-edge identity
                                  :extensions [route-cost]})))
;; => -2
\endtt

Topological dynamic programming handles negative weights because the DAG
has no route back to an earlier node. A~general graph with negative cycles
needs a~different algorithm.

\secc Relax negative edges and detect negative cycles

\noindent{\bf Problem.} A~directed graph may have negative edge weights.
Find the least distance from \code{:s} to each reachable node, or report
whether a~reachable negative cycle makes the distances unbounded below.

\smallskip
\noindent{\bf Solution.} Bellman--Ford repeatedly relaxes every edge:
if reaching \code{from} costs \code{d}, then reaching \code{to} through
that edge costs at most \code{d + weight}. After at most $|V|-1$ passes,
shortest paths without repeated vertices have been considered. If another
pass could still improve a~distance, a~reachable negative cycle exists.
These are synchronous algorithm rounds, expressed directly with
\code{iterate}; they are not fixed-point \code{propagate} updates:

\begtt \hisyntax{clojure}
(defn bellman-ford [nodes edges source]
  (let [relax-edge
        (fn [distances [from to weight]]
          (if-some [from-distance (distances from)]
            (update distances to
                    #(if (some? %)
                       (min % (+ from-distance weight))
                       (+ from-distance weight)))
            distances))
        relax (fn [distances] (reduce relax-edge distances edges))
        distances (nth (iterate relax {source 0})
                       (dec (count nodes)))
        negative-cycle?
        (some (fn [[from to weight]]
                (when-some [from-distance (distances from)]
                  (when-some [to-distance (distances to)]
                    (< (+ from-distance weight) to-distance))))
              edges)]
    (if negative-cycle? :negative-cycle distances)))

(def nodes #{:s :a :b :t})
(def edges [[:s :a 4] [:s :b 3] [:a :b -2]
            [:a :t 5] [:b :t 1]])

[(bellman-ford nodes edges :s)
 (bellman-ford #{:s :a :b}
               [[:s :a 1] [:a :b -2] [:b :a 1]]
               :s)]
;; => [{:s 0, :a 4, :b 2, :t 3} :negative-cycle]
\endtt

The vertex-count bound is essential: with no negative cycle, a~shortest
route needs no repeated vertex. On a~reachable negative cycle, adding one
more edge can keep lowering the distance. These deterministic relaxation
passes are distinct from change-driven fixed-point propagation.

\secc Grow a~minimum spanning tree

\noindent{\bf Problem.} Connect four towns using the least total road length.
The undirected roads have lengths \code{:a-:b} (1), \code{:a-:c} (4),
\code{:b-:c} (2), \code{:b-:d} (5), and \code{:c-:d} (1). Which roads
belong to a~minimum spanning tree, and what is their total length?

\smallskip
\noindent{\bf Solution.} Prim's algorithm repeatedly chooses the lightest
edge leaving the already reached towns. Give each visit the weight of the
edge that joins it, and use that value as best-first's priority. This is
not Dijkstra's algorithm: the priority is the joining edge, not the
cumulative distance from the start:

\begtt \hisyntax{clojure}
(def roads
  {:a [{:state :b :weight 1} {:state :c :weight 4}]
   :b [{:state :a :weight 1} {:state :c :weight 2}
       {:state :d :weight 5}]
   :c [{:state :a :weight 4} {:state :b :weight 2}
       {:state :d :weight 1}]
   :d [{:state :b :weight 5} {:state :c :weight 1}]})

(def joining-edge
  {:start (fn [_] {:joining-edge 0})
   :next (fn [parent transition _]
           {:joining-edge (:weight transition)
            :road [(:state parent) (:state transition)
                   (:weight transition)]})})

(let [tree (search/best-first roads [:a]
                              {:decode-edge identity
                               :extensions [joining-edge]
                               :priority :joining-edge})]
  [(mapv :road (rest tree))
   (transduce (map :joining-edge) + tree)])
;; => [[[:a :b 1] [:b :c 2] [:c :d 1]] 4]
\endtt

Each newly reached town is attached through its cheapest available edge.
The three selected roads have total length four.

\secc Find a~route minimizing its largest climb

\noindent{\bf Problem.} A~hiker crosses this height grid from the upper-left
corner to the lower-right. The effort of a~route is its largest height
difference between consecutive cells. Find a~route with the smallest
possible effort:

\begtt \hisyntax{clojure}
(def heights [[1 2 2]
              [3 8 2]
              [5 3 5]])
\endtt

\smallskip
\noindent{\bf Solution.} Unlike ordinary route length, the effort is not
the sum of the climbs. Accumulate the largest climb seen so far and let
best-first prioritize that value:

\begtt \hisyntax{clojure}
(defn grid-climbs [[row col]]
  (for [[next-row next-col]
        [[(dec row) col] [(inc row) col]
         [row (dec col)] [row (inc col)]]
        :let [here (get-in heights [row col])
              there (get-in heights [next-row next-col])]
        :when there]
    {:state [next-row next-col]
     :climb (abs (- there here))}))

(def effort-extension
  (search/accumulate :effort (constantly 0)
                     (fn [effort transition]
                       (max effort (:climb transition)))))

(->> (search/best-first grid-climbs [[0 0]]
                        {:decode-edge identity
                         :extensions [effort-extension]
                         :priority :effort})
     (filter #(= [2 2] (:state %)))
     first
     (select-keys [:state :effort]))
;; => {:state [2 2], :effort 2}
\endtt

Along any route, adding a~climb can never lower its maximum, so this
priority orders visits by the best effort found so far.

\secc Two travelers meet as quickly as possible

\noindent{\bf Problem.} Two travelers start in \code{:west} and
\code{:east}. On each turn, both must move along one road. A~turn lasts as
long as the slower traveler takes to reach the next city; the faster traveler
waits. Find the minimum elapsed time until they are in the same city.

\smallskip
\noindent{\bf Solution.} A~state is the pair of current cities. Each
successor is one pair of simultaneous moves, and its cost is the longer of
the two road times. Dijkstra search minimizes the sum of these per-turn
durations. This models the joint state directly rather than searching for
each traveler independently:

\begtt \hisyntax{clojure}
(def meeting-roads
  {:west [[:ridge 2]]
   :ridge [[:west 2] [:town 2] [:bridge 5]]
   :east [[:market 1]]
   :market [[:east 1] [:town 2] [:bridge 4]]
   :town [[:ridge 2] [:market 2]]
   :bridge [[:ridge 5] [:market 4]]})

(defn simultaneous-moves [[city-a city-b]]
  (for [[next-a time-a] (meeting-roads city-a)
        [next-b time-b] (meeting-roads city-b)]
    {:state [next-a next-b]
     :weight (max time-a time-b)}))

(->> (search/best-first simultaneous-moves
                        [[:west :east]]
                        {:decode-edge identity
                         :extensions [search/path-extension]})
     (filter (fn [{[city-a city-b] :state}]
               (= city-a city-b)))
     first
     (select-keys [:path :distance]))
;; => {:path [[:west :east] [:ridge :market] [:town :town]]
;;     :distance 4}
\endtt

The two road times are not added within a~turn: one traveler waits for the
other. Modeling the pair as one state and the elapsed turn duration as one
edge cost makes the ordinary weighted search objective match the problem.

\secc Make a~short presentation covering required topics

\noindent{\bf Problem.} A~short presentation must cover \code{:welcome},
\code{:robots}, and \code{:navigation}. Available clips have a~duration and
cover one or more topics. Find a~presentation with minimum total duration.
This is a~small set-cover search.

\smallskip
\noindent{\bf Solution.} A~state records the topics still to cover and the
clips already selected. At each step choose the first remaining topic and
consider clips that cover it. Each transition's weight is the clip duration,
so best-first returns the least-duration complete presentation:

\begtt \hisyntax{clojure}
(def required-topics #{:welcome :robots :navigation})
(def clips
  [{:id :short-welcome :duration 2 :topics #{:welcome}}
   {:id :robot-navigation :duration 5 :topics #{:robots :navigation}}
   {:id :robot-basics :duration 4 :topics #{:robots}}
   {:id :navigation-demo :duration 2 :topics #{:navigation}}
   {:id :full-introduction :duration 8 :topics #{:welcome :robots}}])

(defn add-clip [{:keys [remaining selected]}]
  (when-let [topic (first (sort remaining))]
    (for [{:keys [id duration topics]} clips
          :when (contains? topics topic)]
      {:state {:remaining (reduce disj remaining topics)
               :selected (conj selected id)}
       :weight duration})))

(def presentation
  (first (filter #(empty? (get-in % [:state :remaining]))
                 (search/best-first
                  add-clip
                  [{:remaining required-topics :selected []}]
                  {:decode-edge identity}))))

(select-keys presentation [:state :distance])
;; => {:state {:remaining #{}
;;             :selected [:robot-navigation :short-welcome]}
;;     :distance 7}
\endtt

The selected topic makes the branching finite and ensures that every chosen
clip covers at least one outstanding requirement. Best-first compares total
duration, rather than clip count; here one multi-topic clip makes the
seven-minute solution shorter than choosing separate clips for each topic.

\secc Catch a~scheduled connection

\noindent{\bf Problem.} A~traveler leaves \code{:home} at 8:00. A~bus
reaches \code{:junction} at 8:10 and another reaches \code{:museum} directly
at 8:25. From the junction, a~bus leaves at 8:15 and reaches the museum at
8:20. Find the earliest arrival. A~connection can be used only if the
traveler reaches its departure stop in time.

\smallskip
\noindent{\bf Solution.} Store departure and arrival times relative to 8:00.
The space receives the full visit and offers only trips that have not
departed yet. Its transition weight is the wait plus travel time, so the
accumulated Dijkstra distance is the arrival time:

\begtt \hisyntax{clojure}
(def timetable
  {:home [{:to :junction :departs 0 :arrives 10}
          {:to :museum :departs 0 :arrives 25}]
   :junction [{:to :museum :departs 15 :arrives 20}]
   :museum []})

(defn catchable-trips [{:keys [state distance]}]
  (for [{:keys [to departs arrives]} (timetable state)
        :when (<= distance departs)]
    [to (- arrives distance)]))

(->> (search/best-first catchable-trips [:home]
                        {:project identity
                         :decode-edge (fn [[state weight]]
                                        {:state state :weight weight})
                         :extensions [search/path-extension]})
     (filter #(= :museum (:state %)))
     first
     (select-keys [:path :distance]))
;; => {:path [:home :junction :museum], :distance 20}
\endtt

The scheduled edges are search-aware: whether a~trip is a~successor depends
on the current arrival time in \code{:distance}. The default
\code{:visit-key} remains \code{:state}; for a~timetable where earlier arrival
always leaves at least as many future options, a~later visit to the same stop
is dominated. A~static stop-to-stop graph cannot directly express which
departures are still catchable; without search-aware generation, the time
would need to be encoded into states or represented with a~time-expanded
graph.

\secc Keep the details of each leg

\noindent{\bf Problem.} A~travel network has these directed trips:
\code{:home -> :museum} by express train (11~hours), \code{:home ->
:village} by ferry (4), \code{:home -> :park} by tram (1),
\code{:village -> :museum} by bus (7), \code{:village -> :hill} by coach
(2), \code{:park -> :museum} by taxi (8), \code{:park -> :village} on
foot (1), and \code{:hill -> :museum} by rail (2). There are no other
trips. Find the least-time route from \code{:home} to \code{:museum}, and
report each leg's mode and time as well as the total time.

\smallskip
\noindent{\bf Solution.} Give each edge its destination, time, and mode.
Decode these into transition maps and use \code{trail-extension} to retain
the full transitions along the selected route:

\begtt \hisyntax{clojure}
(def labelled-travel
  {:home [[:museum 11 :express]
          [:village 4 :ferry]
          [:park 1 :tram]]
   :village [[:museum 7 :bus]
             [:hill 2 :coach]]
   :park [[:museum 8 :taxi]
          [:village 1 :foot]]
   :hill [[:museum 2 :rail]]
   :museum []})

(->> (search/best-first labelled-travel [:home]
                        {:decode-edge (fn [[state weight mode]]
                                        {:state state
                                         :weight weight
                                         :mode mode})
                         :extensions [search/trail-extension]})
     (filter #(= :museum (:state %)))
     first
     (select-keys [:trail :distance]))
;; => {:trail [{:state :home}
;;             {:state :park, :weight 1, :mode :tram}
;;             {:state :village, :weight 1, :mode :foot}
;;             {:state :hill, :weight 2, :mode :coach}
;;             {:state :museum, :weight 2, :mode :rail}],
;;     :distance 6}
\endtt

The trail preserves edge data that a~state-only path would lose.

\secc Find a~route matching a~label pattern

\noindent{\bf Problem.} A~directed network has labeled edges. Find the
shortest route from \code{:start} to \code{:target} whose edge labels match
one or more repetitions of the pattern \code{ab}.

\smallskip
\noindent{\bf Solution.} Carry both the current graph node and the state
of a~small finite automaton in the search state. The automaton accepts after
a~\code{b} follows an \code{a}; from that accepting state, another \code{a}
can begin the next repetition. Thus ordinary BFS explores the product of the
graph and the pattern recognizer, while the default visit identity keeps
distinct automaton states at the same graph node separate:

\begtt \hisyntax{clojure}
(def labelled-graph
  {:start [{:to :middle :label "a"}
           {:to :wrong :label "b"}]
   :middle [{:to :target :label "b"}]
   :target [{:to :middle :label "a"}]
   :wrong [{:to :target :label "a"}]})

(defn matching-edges [[node phase]]
  (for [{:keys [to label]} (labelled-graph node)
        :let [next-phase
              (case [phase label]
                [:need-a "a"] :need-b
                [:need-b "b"] :accept
                [:accept "a"] :need-b
                nil)]
        :when next-phase]
    {:state [to next-phase] :label label}))

(let [route (->> (search/bfs matching-edges [[:start :need-a]]
                             {:decode-edge identity
                              :extensions [search/trail-extension]})
                (filter #(= [:target :accept] (:state %)))
                first)]
  (mapv :label (rest (:trail route))))
;; => ["a" "b"]
\endtt

This expresses a~regular constraint on one route. It is not a~declarative
graph-pattern query: matching patterns that bind several nodes and
relationships together or impose subgraph-matching semantics may require a~
different abstraction.

\secc Guide best-first search with an estimate

\noindent{\bf Problem.} A~delivery network has these directed routes:
\code{:home -> :hill} takes 1~hour, \code{:home -> :river} takes 2,
\code{:home -> :market} takes 4, \code{:hill -> :goal} takes 5,
\code{:river -> :goal} takes 2, and \code{:market -> :goal} takes 2;
there are no other routes. The estimated remaining times are 3~hours from
home, 4 from the hill, 2 from the river, 2 from the market, and 0 from the
goal. Find the least-time route.

\smallskip
\noindent{\bf Solution.} Best-first can prioritize each visit by the time
already travelled plus the estimated time remaining. This is A* search; the
estimates here are consistent, so the first settled goal has minimum travel
time:

\begtt \hisyntax{clojure}
(def delivery-routes
  {:home [[:hill 1] [:river 2] [:market 4]]
   :hill [[:goal 5]]
   :river [[:goal 2]]
   :market [[:goal 2]]
   :goal []})
(def remaining-time
  {:home 3 :hill 4 :river 2 :market 2 :goal 0})

(->> (search/best-first delivery-routes [:home]
                        {:decode-edge (fn [[state weight]]
                                        {:state state :weight weight})
                         :extensions [search/path-extension
                                      search/distance-extension]
                         :priority (fn [{:keys [state distance]}]
                                     (+ distance
                                        (remaining-time state)))})
     (filter #(= :goal (:state %)))
     first
     (select-keys [:path :distance]))
;; => {:path [:home :river :goal], :distance 4}
\endtt

The estimate guides the order of exploration; it does not change the route
cost. For this once-per-state best-first search, use a~consistent estimate
that is zero at the goal.

\secc Reach the goal with limited fuel

\noindent{\bf Problem.} A~vehicle starts at \code{:home} with 5~units of
fuel and cannot refuel. Its only directed routes, annotated with fuel used,
are
\code{:home -> :junction} (4), \code{:home -> :detour} (1),
\code{:detour -> :junction} (1), and \code{:junction -> :goal} (3).
Can it reach \code{:goal} without any move using more fuel than remains?
Arriving with zero fuel is allowed.

\smallskip
\noindent{\bf Solution.} Fuel remaining affects which routes are possible,
so the space function needs the full visit, not only its location. Track
fuel with an extension and include it in \code{:visit-key}: arriving at the
same junction with different fuel is a~different search position.

\begtt \hisyntax{clojure}
(def fuel-routes
  {:home [[:junction 4] [:detour 1]]
   :detour [[:junction 1]]
   :junction [[:goal 3]]
   :goal []})

(def fuel-left
  (search/accumulate :fuel-left (constantly 5)
                     (fn [fuel transition]
                       (- fuel (:fuel-used transition)))))

(defn affordable-routes [{:keys [state fuel-left]}]
  (filter (fn [[_ fuel-used]] (<= fuel-used fuel-left))
          (fuel-routes state)))

(->> (search/bfs affordable-routes [:home]
                 {:project identity
                  :visit-key (juxt :state :fuel-left)
                  :decode-edge (fn [[state fuel-used]]
                                 {:state state :fuel-used fuel-used})
                  :extensions [fuel-left search/path-extension]})
     (filter #(= :goal (:state %)))
     first
     (select-keys [:path :fuel-left]))
;; => {:path [:home :detour :junction :goal], :fuel-left 0}
\endtt

The direct arrival at \code{:junction} leaves too little fuel to continue.
If visits were identified only by location, BFS would discard the later
arrival with enough fuel. Including \code{:fuel-left} in \code{:visit-key}
preserves both possibilities.

\secc Collect samples and return to base

\noindent{\bf Problem.} A~small rover starts at \code{[0 0]} with four
units of fuel. It must collect samples at \code{[2 0]} and \code{[2 2]},
refilling to four units at \code{[1 1]}, then return to its starting
position. Each grid move and each refuelling action takes one step. What is
the fewest steps?

\smallskip
\noindent{\bf Solution.} Position alone is not a~complete state: fuel and
collected samples determine which actions remain possible and whether the
goal has been reached. Store all three in each state. The default visit
identity then distinguishes configurations that share a~position but differ
in fuel or samples:

\begtt \hisyntax{clojure}
(def samples {[2 0] :amber, [2 2] :cobalt})
(def station [1 1])
(def base [0 0])
(def capacity 4)

(defn rover-moves [{:keys [position fuel collected]}]
  (let [moves
        (for [[dx dy] [[-1 0] [1 0] [0 -1] [0 1]]
              :let [next [(+ (first position) dx)
                          (+ (second position) dy)]]
              :when (and (every? #(<= 0 % 2) next) (pos? fuel))]
          {:state {:position next
                   :fuel (dec fuel)
                   :collected (cond-> collected
                                (samples next) (conj (samples next)))}})]
    (if (and (= position station) (< fuel capacity))
      (conj (vec moves)
            {:state {:position position
                     :fuel capacity
                     :collected collected}})
      moves)))

(def initial-rover
  {:position base :fuel capacity :collected #{}})

(->> (search/bfs rover-moves [initial-rover]
                 {:decode-edge identity})
     (filter #(and (= base (get-in % [:state :position]))
                   (= #{:amber :cobalt}
                      (get-in % [:state :collected]))))
     first
     :depth)
;; => 12
\endtt

BFS finds the minimum number of actions because moves and refuelling each
count as one step. This is a~product state space: a~physical position is
combined with resource and collection progress.

\secc Build tasks in dependency order

\noindent{\bf Problem.} An application has four tasks: \code{:app} needs
\code{:lib} and \code{:config}; each of those needs \code{:base}; and
\code{:base} has no prerequisites. Find an order in which to perform all
four tasks, with each prerequisite completed before the task that needs it.

\smallskip
\noindent{\bf Solution.} Draw an edge from each task to its prerequisites.
Post-order traversal returns each task after its prerequisites. The shared
\code{:base} prerequisite appears only once:

\begtt \hisyntax{clojure}
(def task-prerequisites
  {:app [:lib :config]
   :lib [:base]
   :config [:base]
   :base []})

(mapv :state
      (search/postorder task-prerequisites [:app]))
;; => [:base :lib :config :app]
\endtt

This ordering is valid because the graph has no dependency cycle. If the
dependency graph has a~cycle, no ordering can put every prerequisite first.

\secc Schedule ordered preparation steps

\noindent{\bf Problem.} A~lunch needs three preparation steps: chop
vegetables (4~minutes), mix dressing (3~minutes), and plate the food
(2~minutes). They use one work surface, so they must be done in that order.
At what minute does each step begin?

\smallskip
\noindent{\bf Solution.} Treat the lunch as depending on its preparation
steps in work order. A~post-order \code{:next} callback sees the completed
earlier steps in the parent's \code{:previous} field, so it can set the next
step's start time to their total duration:

\begtt \hisyntax{clojure}
(def prep-minutes
  {:chop-vegetables 4
   :mix-dressing 3
   :plate-food 2})

(def prep-schedule
  {:next (fn [parent _transition _visit]
           {:starts-at (reduce + 0
                               (map #(prep-minutes (:state %))
                                    (:previous parent)))})})

(->> (search/postorder {:lunch [:chop-vegetables
                                  :mix-dressing
                                  :plate-food]
                          :chop-vegetables []
                          :mix-dressing []
                          :plate-food []}
                         [:lunch]
                         {:extensions [prep-schedule]})
     (filter #(contains? % :starts-at))
     (mapv (juxt :state :starts-at)))
;; => [[:chop-vegetables 0] [:mix-dressing 4] [:plate-food 7]]
\endtt

The callback uses results from earlier siblings to initialize each next
visit. The order matters here because the steps share one work surface.

\secc Restore from the first available copy

\noindent{\bf Problem.} A~document can be restored from a~local copy or,
if that is missing, from a~slower network copy. Check the local copy first.
Which source is used, and can the network lookup be skipped when the local
copy is available?

\smallskip
\noindent{\bf Solution.} Give the preferred source first in the successor
order. Each source reports whether it is available; \code{:more?} lets the
restore visit stop requesting successors once a~completed source succeeds:

\begtt \hisyntax{clojure}
(def available-copies #{:local-copy})

(def restore-extension
  {:more? (fn [{:keys [state]} previous]
            (or (not= state :restore)
                (not-any? :available? previous)))
   :exit (fn [{:keys [state]} successors]
           (if (= state :restore)
             {:source (:state (some #(when (:available? %) %) successors))}
             {:available? (contains? available-copies state)}))})

(mapv #(select-keys % [:state :available? :source])
      (search/postorder {:restore [:local-copy :network-copy]
                           :local-copy []
                           :network-copy []}
                          [:restore]
                          {:extensions [restore-extension]}))
;; => [{:state :local-copy, :available? true}
;;     {:state :restore, :source :local-copy}]
\endtt

The network-copy visit is absent because the local copy succeeded. If it
were unavailable, traversal would continue to the network copy.

\secc Count tasks after their prerequisites

\noindent{\bf Problem.} A~build plan has four tasks: \code{:app} depends on
\code{:compile} and \code{:docs}; both \code{:compile} and \code{:docs}
depend on \code{:parse}; and \code{:parse} has no prerequisites. For each
task, count the distinct tasks reachable by following its prerequisite
links, including the task itself. The shared \code{:parse} task should only
be counted once for \code{:app}.

\smallskip
\noindent{\bf Solution.} A~post-order \code{:exit} callback runs after a~
task's prerequisites have been processed. Carry the set of prerequisite
tasks already counted, union the sets from completed successors, then count
the resulting set:

\begtt \hisyntax{clojure}
(def task-count
  {:exit (fn [{:keys [state]} prerequisites]
           (let [included (into #{}
                                (mapcat #(conj (:included %) (:state %)))
                                prerequisites)]
             {:included (conj included state)
              :size (inc (count included))}))})

(mapv (juxt :state :size)
      (search/postorder {:app [:compile :docs]
                           :compile [:parse]
                           :docs [:parse]
                           :parse []}
                          [:app]
                          {:extensions [task-count]}))
;; => [[:parse 1] [:compile 2] [:docs 2] [:app 4]]
\endtt

This bottom-up computation handles a~shared prerequisite without counting
it twice in the application's total.

\secc Count the shortest routes

\noindent{\bf Problem.} Consider a~directed graph with vertices
\code{:a}, \code{:b}, \code{:c}, \code{:d}, \code{:e}, \code{:f},
\code{:g}, and \code{:h}. Its outgoing edges are \code{:a} to \code{:b},
\code{:c}, and \code{:f}; \code{:b} to \code{:d}; \code{:c} to
\code{:d}; \code{:d} to \code{:e}; \code{:f} to \code{:g}; \code{:g} to
\code{:h}; and \code{:h} to \code{:e}. Every edge has cost one. How many
different routes with the fewest edges lead from \code{:a} to \code{:e}?

\smallskip
\noindent{\bf Solution.} First use BFS to find each state's minimum depth.
Keep only edges that advance exactly one depth level; these edges form a~
directed acyclic graph containing all shortest routes. A~post-order
\code{:exit} callback counts the routes from each state to the goal:

\begtt \hisyntax{clojure}
(def route-graph
  {:a [:b :c :f]
   :b [:d]
   :c [:d]
   :d [:e]
   :f [:g]
   :g [:h]
   :h [:e]
   :e []})

(def depth
  (into {} (map (juxt :state :depth))
        (search/bfs route-graph [:a]
                    {:extensions [search/depth-extension]})))

(defn shortest-edges [state]
  (filter #(= (inc (depth state)) (depth %))
          (route-graph state)))

(->> (search/postorder shortest-edges [:a]
                         {:extensions
                          [{:exit (fn [{:keys [state]} successors]
                                    {:count (if (= state :e)
                                              1
                                              (transduce (map :count) + 0
                                                         successors))})}]})
     last
     :count)
;; => 2
\endtt

The two shortest routes share their final edge from \code{:d} to
\code{:e}; the bottom-up count includes both ways of reaching \code{:d}.

\secc Count ways to break a~word into dictionary words

\noindent{\bf Problem.} Split \code{"catsanddog"} into words from
\code{\{"cat", "cats", "and", "sand", "dog"\}}. How many different
segmentations are possible?

\smallskip
\noindent{\bf Solution.} Use the character offset as a~state. Each
dictionary word that starts at that offset leads to the offset just after
it, so every transition moves forward. A~post-order \code{:exit} callback
can count the segmentations from an offset using the already-counted
successors:

\begtt \hisyntax{clojure}
(def text "catsanddog")
(def dictionary #{"cat" "cats" "and" "sand" "dog"})

(defn next-cuts [start]
  (for [end (range (inc start) (inc (count text)))
        :when (contains? dictionary (subs text start end))]
    end))

(def count-ways
  {:exit (fn [{start :state} suffixes]
           {:ways (if (= start (count text))
                    1
                    (transduce (map :ways) + 0 suffixes))})})

(:ways (last (search/postorder next-cuts [0]
                                 {:extensions [count-ways]})))
;; => 2
\endtt

Each split point is counted once, while distinct choices of the first word
contribute separate ways to the total.

\secc Find the value of tic-tac-toe under perfect play

\noindent{\bf Problem.} Two players take turns marking an empty square on a~
three-by-three board, with \code{:x} moving first. The first player to
complete a~row, column, or diagonal wins; if the board fills first, the game
is a~draw. Assuming both players choose optimally, what is the game's value
for \code{:x}, where a~win is 1, a~loss is -1, and a~draw is 0?

\smallskip
\noindent{\bf Solution.} Each board is a~state and each legal move adds one
mark, so the game graph is acyclic. An \code{:exit} callback evaluates each
board after its children: \code{:x} chooses the largest child value and
\code{:o} chooses the smallest. Terminal boards receive their value directly:

\begtt \hisyntax{clojure}
(def winning-lines [[0 1 2] [3 4 5] [6 7 8]
                    [0 3 6] [1 4 7] [2 5 8]
                    [0 4 8] [2 4 6]])

(defn board-winner [board]
  (some (fn [[a b c]]
          (when (and (board a) (= (board a) (board b) (board c)))
            (board a)))
        winning-lines))

(defn legal-moves [board]
  (when-not (board-winner board)
    (let [mark (if (= (count (filter #{:x} board))
                      (count (filter #{:o} board)))
                 :x :o)]
      (for [i (range 9) :when (nil? (board i))]
        (assoc board i mark)))))

(def game-value
  {:exit (fn [{board :state} children]
           (let [winner (board-winner board)]
             {:value (cond
                       (= :x winner) 1
                       (= :o winner) -1
                       (empty? children) 0
                       (even? (count (remove nil? board)))
                       (apply max (map :value children))
                       :else (apply min (map :value children)))}))})

(let [positions (search/postorder legal-moves
                                    [(vec (repeat 9 nil))]
                                    {:extensions [game-value]})]
  [(count positions) (:value (last positions))])
;; => [5478 0]
\endtt

The empty board's value is zero: optimal play from both sides leads to a~
draw. The same bottom-up traversal used for dependency computations can
therefore evaluate a~game tree by minimax.

\secc Find the edge that closes a~dependency cycle

\noindent{\bf Problem.} A~deployment has two tasks, \code{:api} and
\code{:database}. The API requires the database, and the database requires
the API. Starting at \code{:api}, report each prerequisite edge and indicate
whether it points back to a~task that is still being explored.

\smallskip
\noindent{\bf Solution.} Post-order traversal terminates even when the
graph has cycles. Its \code{:exit} callback receives each successor's
\code{decoded} incoming edge and a~\code{:back-edge} flag, which identifies an
edge to a~task still being explored:

\begtt \hisyntax{clojure}
(def edge-report
  {:exit (fn [_visit successors]
           {:edges (mapv (fn [{:keys [state via back-edge]}]
                           [state (:kind via) back-edge])
                         successors)})})

(mapv (juxt :state :edges)
      (search/postorder {:api [[:database :requires]]
                           :database [[:api :requires]]}
                          [:api]
                          {:decode-edge (fn [[state kind]]
                                          {:state state :kind kind})
                           :extensions [edge-report]}))
;; => [[:database [[:api :requires true]]]
;;     [:api [[:database :requires false]]]]
\endtt

The edge from \code{:database} back to \code{:api} closes the cycle. The
returned order is a~depth-first post-order, not a~valid prerequisite order
when the graph contains a~cycle.

\sec Technical documentation

This section explains the implementation in the order a~search uses it:
visits and transitions, lazy frontiers, shared traversal, then the algorithms
that need their own traversal rules. The code excerpts are taken directly from the runnable source files; those
files remain the complete source of truth. The embedded Lua below locates
named top-level Clojure forms and finds the current excerpt line numbers
whenever OpTeX builds this document.

\_newcount\sourcefirstline
\_newcount\sourcelastline
\directlua{
  local excerpt_bounds = {
    visits = {"(ns state-space.core", "(defn- with-required"},
    ["preparing-search"] = {"(defn- with-required", "(defn- explore"},
    ["shared-traversal"] = {"(defn- explore", "(defn postorder"},
    postorder = {"(defn postorder", "(defn propagate"},
    propagation = {"(defn propagate", "(defn- enqueue"},
    ["best-first"] = {"(defn- enqueue"}
  }

  function state_space_excerpt_range(path, name)
    local file = assert(io.open(path, "r"))
    local lines, line_count = {}, 0
    for line in file:lines() do
      line_count = line_count + 1
      lines[line_count] = line
    end
    file:close()
    local bounds = assert(excerpt_bounds[name],
                         "unknown source excerpt: " .. name)
    local function find_line(prefix)
      if not prefix then return nil end
      local found
      for line_number, line in ipairs(lines) do
        if line:sub(1, string.len(prefix)) == prefix then
          assert(not found, "duplicate source excerpt boundary: " .. prefix)
          found = line_number
        end
      end
      assert(found, "missing source excerpt boundary: " .. prefix)
      return found
    end
    local first_line = find_line(bounds[1])
    local next_line = find_line(bounds[2])
    local last_line = next_line and next_line - 1 or line_count
    while last_line >= first_line and lines[last_line] == "" do
      last_line = last_line - 1
    end
    assert(first_line and last_line and first_line <= last_line,
           "missing or empty source excerpt markers: " .. name)
    local backslash = string.char(92)
    tex.sprint(backslash .. "sourcefirstline=" .. first_line ..
               backslash .. "sourcelastline=" .. last_line)
  end
}

\load[doc]
\hyperlinks \Blue \Green
\def\bgverbcolor{\setcmykcolor{0 0 .12 0}}

\secc Visits, transitions, and extensions

The algorithms exchange visit maps. \code{:state} is the problem's state;
extensions add derived values to the same map. A~raw successor is either
already a~state or is decoded into a~transition map. Its \code{:state}
becomes the child's state, while the remaining transition fields are passed
to \code{:next} callbacks. This gives path, depth, trail, and application
specific extensions one shared lifecycle.

\begingroup
\typosize[8/9]
\ttline=-2
\directlua{state_space_excerpt_range("src/state_space/core.clj", "visits")}
\verbinput \hisyntax{clojure}\vitt{src/state_space/core.clj} (\the\sourcefirstline-\the\sourcelastline) src/state_space/core.clj
\endgroup

\secc Preparing a search

Options are resolved once at the start of an algorithm. The space is called
with the projected visit, and \code{:expand?} can make a visit a leaf without
removing it from the results. Decoding then gives each raw successor a
uniform transition-map shape. The local \code{map-lazy} helper delays its
mapping function one item at a time; it is used where preparing each start
visit or decoded transition should wait until that item is requested.

\begingroup
\typosize[8/9]
\ttline=-2
\directlua{state_space_excerpt_range("src/state_space/core.clj", "preparing-search")}
\verbinput \hisyntax{clojure}\vitt{src/state_space/core.clj} (\the\sourcefirstline-\the\sourcelastline) src/state_space/core.clj
\endgroup

\secc Lazy work and frontiers

The results of the searches are lazy, but laziness here does not promise that
exactly one callback or successor is processed for each requested result.
\code{map-lazy} invokes its mapping function one item at a time; it is used
where start visits and post-order transitions are prepared. Other sequence
operations, such as \code{keep} in the shared traversal, can realize a chunk
when their input is chunked. The space itself also controls how its
successors are produced, so consuming one raw successor may realize several
values from a chunked input. Best-first has an additional requirement: it
must examine all successors of the expanded visit before it can choose the
next visit by priority.
Propagation updates are processed one state at a time, but all of that
state's successors are joined before its progress record is returned.

DFS and BFS keep pending work in a frontier of lazy chunks. A stack takes the
newest chunk first; a queue takes the oldest. Both consume items from the
current chunk in order, so successors retain their order without first being
flattened into one collection.

\begingroup
\typosize[8/9]
\ttline=-2
\verbinput \hisyntax{clojure}\vitt{src/state_space/internal/frontier.clj} (1-) src/state_space/internal/frontier.clj
\endgroup

\secc Shared traversal: DFS, BFS, and simple paths

\code{explore} owns the mechanics common to depth-first and breadth-first
search: pop a visit, optionally deduplicate it, yield it, then lazily prepare
its children and push them onto the frontier. Supplying a stack or queue
changes the order without duplicating this traversal logic. \code{simple-paths}
uses the same depth-first engine but replaces global deduplication with a
per-route refinement. The route's identity keys live in metadata, so they
control pruning without becoming part of the public visit.

\begingroup
\typosize[8/9]
\ttline=-2
\directlua{state_space_excerpt_range("src/state_space/core.clj", "shared-traversal")}
\verbinput \hisyntax{clojure}\vitt{src/state_space/core.clj} (\the\sourcefirstline-\the\sourcelastline) src/state_space/core.clj
\endgroup

\secc Post-order traversal

Post-order traversal cannot use the shared loop: a visit must be returned
only after its successors have finished. Its explicit pending items represent
entering a visit and resuming its parent. The local operations make those
steps explicit: \code{enter} marks a visit known and schedules its
continuation; \code{continue} records a completed child and schedules the
next one; \code{finish} runs bottom-up callbacks and emits the visit.
\code{known} stores the first visit for each key; \code{done} distinguishes
completed nodes from nodes still active on the current path. When a child
finishes, the parent records it in \code{:previous}; that same vector is
passed to \code{:next}, \code{:more?}, and ultimately \code{:exit}. Checking
\code{:more?} before consuming the next transition lets the algorithm stop
before decoding or visiting a remaining successor. A chunked source may
nevertheless realize multiple raw successors as one is consumed.

\begingroup
\typosize[8/9]
\ttline=-2
\directlua{state_space_excerpt_range("src/state_space/core.clj", "postorder")}
\verbinput \hisyntax{clojure}\vitt{src/state_space/core.clj} (\the\sourcefirstline-\the\sourcelastline) src/state_space/core.clj
\endgroup

\secc Best-first traversal

Best-first orders pending visits by priority, so it uses a sorted map rather
than the DFS/BFS frontier. \code{best} records the best priority queued for
each visit key; when a better route is found, the older queue entry is left
in place and discarded as stale when popped. \code{settled} ensures each key
is emitted once. A monotonically increasing sequence number makes equal
priorities stable in insertion order. The default cumulative-distance
priority has Dijkstra's shortest-path behavior when edge weights are
non-negative.

\begingroup
\typosize[8/9]
\ttline=-2
\directlua{state_space_excerpt_range("src/state_space/core.clj", "best-first")}
\verbinput \hisyntax{clojure}\vitt{src/state_space/core.clj} (\the\sourcefirstline-\the\sourcelastline) src/state_space/core.clj
\endgroup

\secc State propagation

Propagation uses the same projected visits and decoded transitions as the
searches, but its worklist is organized around accumulated values per state.
A state is requeued when joining a contribution changes its value; pending
states are coalesced so a state already waiting in the FIFO queue is not
queued twice. Each progress record contains a persistent snapshot of the value
map, the processed state, and its changed destinations. The initial snapshot
is yielded before any state is expanded.

\begingroup
\typosize[8/9]
\ttline=-2
\directlua{state_space_excerpt_range("src/state_space/core.clj", "propagation")}
\verbinput \hisyntax{clojure}\vitt{src/state_space/core.clj} (\the\sourcefirstline-\the\sourcelastline) src/state_space/core.clj
\endgroup

\sec Acknowledgments and further reading

Developed with help from an AI coding companion, who never once complained
about the parentheses.~\emoji{smile}

The library's design and examples draw inspiration from the projects and
works listed below.

\begitems
* \ulink[https://nographs.readthedocs.io/en/latest/]{NoGraphs documentation},
  for on-the-fly graph analysis and a broad set of traversal examples.
* \ulink[https://github.com/clj-commons/loom]{Loom}, a Clojure graph library,
  as a comparison for graph representations and algorithms in this language.
* The Rust \ulink[https://docs.rs/pathfinding/latest/pathfinding/]{pathfinding
  crate}, as a comparison for generic pathfinding and graph algorithms.
* Chris Hanson and Gerald Jay Sussman,
  \ulink[https://doi.org/10.7551/mitpress/12821.001.0001]
  {{\it Software Design for Flexibility: How to Avoid Programming Yourself
  into a Corner}}, for its flexible and composable design perspective.
* David~L. Poole and Alan~K. Mackworth,
  \ulink[https://artint.info/3e/html/ArtInt3e.html]
  {{\it Artificial Intelligence: Foundations of Computational Agents},
  3rd~edition}, for its treatment of state-space search, planning, and
  problem representation.
* Stuart Russell and Peter Norvig, {\it Artificial Intelligence: A Modern
  Approach}, 3rd~edition (2010), and the
  \ulink[https://aimacode.github.io/aima-exercises/]{AIMA exercises}, for
  problems and concepts used to examine the scope of the search model.
* The \ulink[https://leetcode.com/problemset/]{LeetCode problem collection},
  for additional problem ideas used to exercise the library's search model.
\enditems

\bye
